% H9-T16: semistable degeneration attack on the fixed-axis gate.
% Run JAC1964-20261121; /root/h9_height; GPT-6 Astra Pro.
% New deductions are separated from cited weight-spectral-sequence inputs.

\subsection*{H9-T16. A semistable axis degeneration and its exact lifting obstruction}

Let $A/k$ be an abelian variety of dimension $2m$, $m\ge2$, over an
algebraically closed field. Let $H_1,H_2\in|L|$ meet transversely,
put $h=c_1(L)$ and $W=H_1\cap H_2$, and let $E/k$ be an elliptic
curve with origin $0$. Thus $\dim W=2m-2$. Write
$\kappa:W\hookrightarrow A$ and
\[
 \mathcal A_W=\ker\bigl(\kappa_*:
  \operatorname{Alg}^{m-1}(W)\longrightarrow H^{2m+2}(A)(m+1)\bigr).
\]
As proved in H9-T7, every element of $\mathcal A_W$ is primitive
for $h_W$. The remaining fixed-axis sign is
$(-1)^{m-1}\deg_W(ww')$ on this space.

\paragraph{Two necessary corrections to the proposed smoothing.}
An ample product line bundle
$\operatorname{pr}_A^*L\otimes\operatorname{pr}_E^*\mathcal O_E(p)$
with $\deg(p)=1$ does not give the intended smoothing. Indeed
$h^0(E,\mathcal O_E(p))=1$, and the Kunneth formula for global
sections makes every section a product. Every divisor in that
complete linear system contains $A\times\{p\}$. There is no
irreducible member in this system.

Instead choose reduced effective divisors
\[
 P_1=p_{11}+\cdots+p_{1q},\qquad
 P_2=p_{21}+\cdots+p_{2q}
\]
of the same degree $q\ge2$, with disjoint supports. For embeddings
one may take $q\ge3$. We may, and for the Chow statement below do,
choose all their points torsion. They can be chosen linearly
 equivalent, for example with sum zero in the elliptic group.
On $A\times E$ put
\[
 \mathcal L_i=L\boxtimes\mathcal O_E(P_i),\qquad
 D_i^0=(H_i\times E)\cup(A\times P_i).
\]
The reduced central complete intersection has components
\[
 Y_0=V_0\cup\bigcup_a V_a,
 \quad V_0=W\times E,
 \quad V_a=H_2\times\{p_{1a}\}\ \text{or}\ H_1\times\{p_{2a}\}.
                                                               \tag{H9.64}
\]
It has no triple intersections. Each double intersection is a copy
$D_a=W$, attached as $W\times\{p_a\}\subset V_0$ and as the
Cartier divisor $W\subset H_i$ in its leaf.

Although (H9.64) is simple normal crossing, it is not the special
fibre of a regular semistable total space with these unchanged
components. Along $D_a$, the product of its two normal bundles is
\[
 N_{D_a/V_0}\otimes N_{D_a/V_a}=\mathcal O_W(h),             \tag{H9.65}
\]
not the trivial bundle. For two components of the special fibre of
a regular strictly semistable scheme with no third component
meeting them, that product is trivial: restrict the principal
special-fibre divisor to their intersection. Equation (H9.65)
therefore supplies a second geometric obstruction to the unmodified
proposal.

\paragraph{A corrected semistable model.}
Take a one-parameter smoothing of the two product sections,
locally written $s_i t_i+t g_i=0$. At a double intersection over a
zero of $t_1$, the other elliptic section is a unit. Eliminate the
second equation. The remaining equation has the form
\[
                         xy=t a(z)+O(t^2),                 \tag{H9.66}
\]
where $a$ is a section of $\mathcal O_W(h)$, up to a nonzero scalar.
Assume its zero divisor $Z_a\subset W$ is smooth, and make the
analogous assumption at every double intersection. Away from these
zero divisors the total space is already regular and semistable.
Near a zero of $a$, an etale change of a transverse coordinate puts
(H9.66) into the form $xy=tz$ times a smooth factor.

Blow up the ideal $(y,t)$ of the corresponding leaf. One chart has
$t=yv$ and $x=vz$, so its total space is smooth and the special
fibre is $yv=0$. The other chart has $y=tu$ and $z=xu$, with
special fibre $t=0$; it is smooth as well. This modification leaves
$V_0=W\times E$ unchanged and replaces the leaf by
\[
                    V_a'=\operatorname{Bl}_{Z_a}H_i.
                                                               \tag{H9.67}
\]
The double divisor remains isomorphic to $W$. Its normal bundle in
the modified leaf is
$\mathcal O_W(h)(-Z_a)=\mathcal O_W$, as required. The exceptional
locus over $Z_a$ is a projective line bundle of dimension one less
than a component; it is not a new irreducible component of the
special fibre. The modifications at the distinct elliptic attaching
points have disjoint centres.

Thus, when the generic complete intersection is smooth and the
indicated $Z_a$ are smooth, these explicit modifications give a
proper strictly semistable model $\mathcal Y/R$ with central
components $V_0,V_a'$ and no triple intersections. These hypotheses
are kept explicit. They can be arranged by general choices in
sufficiently ample systems. A geometric generic complete intersection
can be approached by a trait through the product-section boundary
point: blow up that point in its parameter space and take the local
ring at the generic point of the exceptional divisor. Its residue
field contains the generic smoothing directions. Extension of the
residue field does not change the intersections under discussion.
No assertion that an arbitrary one-parameter smoothing automatically
satisfies the terminal-pencil hypothesis of H9-T13 is being made.

\paragraph{(i) The surviving cohomological block.}
Put $n=2m-1$ and let $D=\coprod_aD_a$ and
$V=V_0\amalg\coprod_aV_a'$. For every invertible prime $\ell$,
the central term of the weight spectral sequence is
\[
 E_2^{0,n}
   =\frac{\ker\{H^n(V)\xrightarrow{\rho}H^n(D)\}}
          {\operatorname{im}\{H^{n-2}(D)(-1)
                                \xrightarrow{\gamma}H^n(V)\}}.
                                                               \tag{H9.68}
\]
Here $\rho$ is alternating restriction and $\gamma$ alternating
Gysin. The weight spectral sequence and these differentials are
T.~Saito, \emph{Weight spectral sequences and independence of
$\ell$}, Corollary 2.2.4 and Proposition 2.2.6, author pp.~21,
23--24. Its degeneration at $E_2$ for an arbitrary henselian
valuation field is T.~Ito, \emph{Weight-monodromy conjecture over
equal characteristic local fields}, Theorem 1.1(1), author p.~2.
The latter theorem explicitly includes arbitrary residue fields.

The Kunneth summand
\[
 M=H^{2m-2}(W)\otimes H^1(E)
       \subset H^n(V_0)                                    \tag{H9.69}
\]
is in the kernel of every restriction, because $H^1(E)$ restricts
to zero at a point. The $V_0$ component of the Gysin map from a
double divisor is
\[
 H^{2m-3}(W)(-1)\longrightarrow H^{2m-3}(W)\otimes H^2(E),
       \qquad u\longmapsto u\otimes[p_a].
\]
It has zero projection to $M$. Consequently the Kunneth projection
onto $M$ splits the complex in (H9.68) at its middle term. In
particular $M$ is a direct summand of $E_2^{0,n}$.

The direct-sum Poincare form on the components restricts on $M$ to
\[
 (w\otimes u,w'\otimes v)
       \longmapsto \left(\int_Www'\right)\left(\int_Euv\right).
                                                               \tag{H9.70}
\]
There is no Koszul sign here: the degrees on $W$ are even. This
form on the full space $M$ is nondegenerate. Its pairings with the
other Kunneth summands of $V_0$ vanish by degree, and the other
components have disjoint direct-sum factors. Restriction and Gysin
are adjoint, so it descends to the indicated nondegenerate summand
in (H9.68). Degeneration at $E_2$ identifies it with a summand of
$\operatorname{Gr}^W_nH^n(Y_{\bar\eta})$. This does not choose a
canonical splitting into the unfiltered nearby cohomology.

Since there are no triple intersections, the only possible weights
in $H^n(Y_{\bar\eta})$ are $n-1,n,n+1$. The monodromy operator
lowers weight by two. It therefore kills $W_n$, and in particular
all invariant lifts of (H9.69). Equivalently its induced map out of
$E_2^{0,n}$ has zero target. After the twist by $m-1$, an algebraic
line $\mathbf Q_\ell[w]\otimes H^1(E)$ is a middle-weight-one
elliptic block, not a toric monodromy block.

These computations prove cohomological survival, with no
Hodge-standard conjecture or algebraic lifting assumption.

\paragraph{(ii) Rational operational-Chow gluing really holds.}
Use a second elliptic copy $E_T$ as parameter and the normalized
Poincare divisor
\[
 P_E=\Delta_E-E_T\times\{0\}-\{0\}\times E
             \quad\text{on }E_T\times E.
\]
For a rational cycle $w\in\operatorname{CH}^{m-1}(W)_{\mathbf Q}$,
put on the smooth components of $E_T\times Y_0$ the tuple
\[
             c_0=w\boxtimes P_E,
       \qquad c_a=0\quad\text{on }E_T\times V_a'.           \tag{H9.71}
\]
Its restriction to a double intersection is
$w\boxtimes\mathcal O_E(p_a-0)$. If every $p_a$ is torsion, this
is zero in rational Chow. Thus all double restrictions agree.

For completeness, equality of these restrictions gives an actual
\emph{operational} Chow class, not just a formal list of classes.
The normalization $\coprod V_i\to Y_0$ is a proper envelope, and
remains an envelope after every base change. The envelope exact
sequence for Chow homology is
\[
 \operatorname{CH}_*(\widetilde T\times_T\widetilde T)
 \xrightarrow{p_{1*}-p_{2*}}\operatorname{CH}_*(\widetilde T)
 \longrightarrow\operatorname{CH}_*(T)\longrightarrow0.
                                                               \tag{H9.72}
\]
This is Stacks Project, Lemma 42.22.4, tag 0GU8. Apply it to an
arbitrary test scheme $T\to E_T\times Y_0$. To act on a Chow
class on $T$, lift it to the component envelope, act there by the
classes (H9.71), and push back. Independence of the chosen lift
follows from (H9.72): on a mixed fibre product the two pulled-back
operations agree by their equality on the corresponding double
intersection; on a same-component fibre product they agree
identically. Compatibility with proper pushforward, flat pullback
and refined Gysin follows from those compatibilities on each
component. This defines the desired operational class.

Thus even the rational Chow gluing condition is satisfied. It says
nothing about surjectivity of restriction from a regular smoothing
or specialization from its generic Chow group.

\paragraph{(iii) A precise obstruction to lifting with the same elliptic parameter.}
Assume now that the generic fibre of the smoothing is, after extension
of fields, the geometric generic member of an ample complete-
intersection family in $A\times E$ for which the terminal pencil
has the following established properties:
\begin{enumerate}
\item the variable system
$\mathbb V_\ell=\ker(k_*:H^n(Y_{\bar\eta})
       \to H^{n+4}(A\times E)(2))$
is absolutely irreducible after restriction to every finite-index
subgroup of its geometric monodromy group;
\item $\dim\mathbb V_\ell>2$;
\item the ambient projector is algebraic, as given by
$(L_{h'}^2)^{-1}k_*$ for the ample complete-intersection class $h'$.
\end{enumerate}
For the genuine terminal even-dimensional Lefschetz pencils in
H9-T13, property (1) is SGA 7 II, XVIII, 6.7.2; property (3) holds
because $A\times E$ is abelian. Property (2) is an explicit rank
condition, satisfied in sufficiently ample examples. It is not
being inferred merely from the word ``nodal.''

Let $0\ne[w]\in\mathcal A_W$. There is \emph{no} rational
algebraic family on $E_T\times Y_{\bar\eta}$ extending to the
semistable model whose componentwise degree-one parameter
realizations are $[w\boxtimes P_E]$ on $V_0$ and zero on the
leaves. More generally, it is enough to prescribe that its entire
middle-weight realization is this summand (H9.69), with zero
remaining component and zero ambient pushforward.

Here is the proof with the two required compatibility steps.
Suppose such a family $\Gamma$ exists. Project to degree one on
$E_T$. The weight spectral sequence for the good-reduction elliptic
parameter has its $H^1$ in column zero only. Functoriality for
algebraic correspondences gives a map to the weight spectral
sequence of $\mathcal Y$. On $E_2^{0,n}$ it is
\[
              H^1(E_T)(1-m)\longrightarrow M,
                       \quad u\longmapsto[w]\otimes u,
                                                               \tag{H9.73}
\]
under the Poincare identification furnished by $P_E$, up to its
fixed nonzero conventional scalar. It is nonzero because $[w]\ne0$.
Both sequences degenerate at $E_2$, so the corresponding generic
cohomological map $f_\Gamma$ is nonzero. This is a use of
T.~Saito, Proposition 2.4.4, author pp.~35--36, in its actual
\emph{generic correspondence to special weight complex} direction;
it does not reverse that proposition to manufacture $\Gamma$.

Next push this correspondence to the constant ambient variety
$E_T\times A\times E$. Proper pushforward and refined restriction
commute. On the central fibre its mixed degree-one class is
\[
             (\kappa_*[w])\boxtimes[P_E]=0,                 \tag{H9.74}
\]
with zero contributions from the leaves. The cohomology sheaf of
the constant ambient family is constant, so its generic class is
also zero. Equivalently this follows from functoriality of the
weight sequences, since both the parameter and ambient target
have good reduction and therefore only the relevant column-zero
terms. Consequently $f_\Gamma$ has image in $\mathbb V_\ell$.

The image has rank at most two. After a finite extension over which
the algebraic family is defined, it is stable under the geometric
monodromy of the generic complete intersection. By property (1) it
must be zero or the whole variable space. The second possibility
contradicts property (2), while the first contradicts (H9.73).
This proves the obstruction.

There is no hidden restriction to cycles with coefficients algebraic
over the function field of the global parameter space. If the
trait is a completion or has enlarged algebraically closed constants,
first spread the cycle over a finitely generated extension of the
algebraic closure of that parameter field. On the constant smooth
projective variety $E_T\times Y$, the relative cycle class is a
constant section. Specializing the spread to a rational point over
that algebraically closed field preserves both its nonzero mixed
class and its zero ambient pushforward. The resulting algebraic
cycle is defined over a finite extension of the original generic
parameter field, to which the finite-index argument applies.
This uses ordinary spreading and smooth proper base change, not
surjectivity of Chow specialization in the semistable family.

\paragraph{A nonzero axis example, so the obstruction is not vacuous.}
Take $A=E_0^4$ and let $a_i$ be the pullback of the origin divisor
from its $i$-th factor. Put $h=d(a_1+a_2+a_3+a_4)$ for $d$ sufficiently
large, and let
$C=E_0\times\{0\}^3$. Two general divisors in $|h|$ containing
$C$ have smooth transverse intersection $W$ for sufficiently large
$d$. Indeed their normal derivatives along $C$ are two general
sections of the rank-three conormal bundle twisted by $h$; the
rank-drop locus has expected codimension two, larger than
$\dim C=1$. Serre generation supplies the necessary simultaneous
evaluations, and Bertini gives smoothness away from $C$.

Define the rational ambient divisor class
\[
 a=\frac{-2a_1+a_2+a_3+a_4}{6d^2}.
\]
Since $a_i^2=0$, direct multiplication gives
$h^2a=a_2a_3a_4=\kappa_*[C]$. Thus
\[
                 w=[C]-\kappa^*a\in\mathcal A_W.
\]
Adjunction on the smooth surface $W$, with $K_W=2h_W$ and
$g(C)=1$, gives $C^2=-2d$. Projection formula gives
\[
 \int_Ww^2
   =-2d-\int_Aa\,a_2a_3a_4
   =-2d+\frac1{3d^2}\ne0.                                 \tag{H9.75}
\]
In particular $[w]\ne0$. This calculation is valid in every
characteristic.

For a complete intersection threefold of two divisors of class
$h'$ in an abelian fivefold, weak Lefschetz, duality and
$c(T_Y)=(1+h')^{-2}|_Y$ give
\[
 \chi(Y)=-4\int h'^5,\qquad
 b_3(Y)=4\int h'^5+72,\qquad
 \dim\mathbb V_\ell=4\int h'^5-48.                          \tag{H9.76}
\]
If the extra elliptic degree is also $d$, then
$\int h'^5=120d^5$, so the variable rank is much larger than two.
Taking $d$ a sufficiently large multiple of three puts $h'$ in a
degree-at-least-two multiple of a product projective embedding,
which supplies the terminal-pencil hypothesis of H9-T13.
This yields explicit nonzero axis classes satisfying (i) and (ii)
that cannot lift in the sense of (iii). It refutes that particular
lifting criterion, not the Hodge-standard sign: for this surface
example the sign is already positive by the surface index theorem.

\paragraph{What the toric pairing can and cannot add.}
The common integral toric monodromy form used in H9-T6 concerns the
weight-zero and weight-two pieces of a semistable abelian variety's
$H^1$. The block (H9.69), after twisting, has weight one and is
annihilated by $N$. Consequently there is no nonzero rank-one
monodromy coefficient on this block to compare with the
Picard--Lefschetz coefficient. The all-prime rational-square
argument of H9-T6 therefore does not determine the sign of
$(-1)^{m-1}\int_Ww^2$ here.

A different possible input would be an actual generic incidence
abelian variety whose Neron model has an abelian special quotient
containing the desired elliptic direction, together with an
incidence-specialization identity identifying its pulled-back
polarization with the form in (H9.70). A special elliptic factor of
a larger abelian degeneration need not lift as an elliptic factor
of its generic fibre. Accordingly (iii) alone does not refute this
larger-factor possibility. Neither the weight cocycle nor the
operational-Chow class constructs that generic abelian variety
or proves that its specialization contains this new direction.

In characteristic zero and sufficiently high terminal degree, the
variable system has Hodge level greater than one (SGA 7 II, XX,
Remark 4.3, printed p.~356). Absolute irreducibility then makes
the entire variable higher-Picard factor zero: a nonzero normalized
family would inject a twist of an abelian $H^1$, of level at most
one, onto the whole variable system. Thus in that situation even
a larger variable incidence factor cannot provide the direction.
This is an additional characteristic-zero obstruction to that
specific generic smoothing route, not an argument against the
classical Hodge--Riemann sign itself.

\paragraph{Exact status and next gate.}
The degeneration proves cohomological survival and rational
operational-Chow gluing, but it does not close the arbitrary
positive-characteristic fixed-axis sign. It also gives a genuine
nonlifting theorem under the stated generic monodromy and rank
hypotheses. A deformation constrained to an incidence locus
carrying the axis cycle has different monodromy; the generic
irreducibility argument need not hold on that locus. Positivity
of its particular primitive incidence form would require a new
argument, not an automatic application of H9-T13 to an arbitrary
special subfamily. The distinction between these two loci is the
next geometric question.

\paragraph{A further test on the support-preserving incidence locus.}
There is a stronger obstruction to the most direct constrained
deformation. Write $w=\sum_i n_i[Z_i]$ and suppose a smooth ample
complete intersection $Y'\subset A\times E$, of two divisors
of class $\lambda=h+q[0]_E$, contains every product $Z_i\times E$.
On $Y'$ put $v=\sum_i n_i[Z_i\times E]$, a cycle of codimension
$m-1$. Its ambient pushforward has class
$\kappa_*[w]\otimes1=0$. Weak Lefschetz identifies
$H^{2m-2}(Y')$ with $H^{2m-2}(A\times E)$, and the Gysin map
on this space is multiplication by $\lambda^2$ under that
identification. It is injective: multiplication by $\lambda^3$
is the hard-Lefschetz isomorphism in this degree. Therefore
$[v]=0$. The family obtained by intersecting $v$ with the
Poincare divisor pulled back from $E_T\times E$ has zero odd
realization and hence zero incidence homomorphism by JAC-P01.
Thus preserving these product supports does not produce the
nonzero elliptic direction in (H9.73), even after restricting to
their incidence locus.

Some such support-preserving loci are empty. For a single product
$Z\times E\subset Y'$, write its below-middle class as the
restriction of $u=u_0+u_1+u_2[0]_E$ from $A\times E$.
The equality $\lambda^2u=\kappa_*[Z]\otimes1$ gives
$h^2u_1=0$ and $h^2u_2+2qh u_0=0$. Hard Lefschetz makes
$h^2$ injective on $H^{2m-3}(A)$ and $h$ injective on
$H^{2m-2}(A)$, so $u_1=0$ and
$u_0=-(2q)^{-1}hu_2$. Necessarily
\[
                     \kappa_*[Z]\in h^3H^{2m-4}(A).
\]
For the explicit curve $C\subset E_0^4$ above,
$\kappa_*[C]=a_2a_3a_4$ is not a multiple of $h^3$.
Consequently no such smooth ample complete intersection contains
$C\times E$. A different cycle correspondence on a constrained
deformation is not excluded by this argument, but its existence
and its primitive incidence sign are new requirements; the
support-preserving construction does not satisfy them.

The weight inputs used above were directly read in the primary
sources
\url{https://www.ms.u-tokyo.ac.jp/~t-saito/pp/wmr2r.pdf} and
\url{https://arxiv.org/pdf/math/0308141}; local copies are
\texttt{work/h9/Saito\_weight2003.pdf} and
\texttt{work/h9/Ito\_weight2003.pdf}.
The envelope exact sequence was directly read at
\url{https://stacks.math.columbia.edu/tag/0GU4}.
The explicit local equations, block calculation, operational descent
and lifting obstruction are deductions proved in this run.
No formal verification or historical originality is claimed.
